
%%%%%%%%%%%%%%%%%%%%%%%%%%%%%%%%%%%%%%%%
%           main body
%%%%%%%%%%%%%%%%%%%%%%%%%%%%%%%%%%%%%%%%

%-----------------------------------
%	TITLE PAGE
%-----------------------------------

\title[第四章\\
勒贝格积分]{\texorpdfstring{\LARGE \bfseries 第四章\quad 勒贝格积分}{第四章 勒贝格积分}} % The short title appears at the bottom of every slide, the full title is only on the title page
\author[\S 4.4\\
R积分与L积分] % Your institution as it will appear on the bottom of every slide, maybe shorthand to save space
{\Large \bfseries \S 4.4 R积分与L积分的比较}
% \author{Singular Value Decomposition} % Your name
% \date{\today} % Date can be changed to a custom date
\date{}

%------------------------------------------------

\begin{document}

%------------------------------------------------

\begin{frame}

\titlepage

\end{frame}

%------------------------------------------------

\section[引言]{引言: 为什么要比较}

%------------------------------------------------

\begin{frame}{引言: Riemann 积分与它的困境}

\begin{itemize}
\item \textbf{Riemann 积分} (数学分析): 对定义域作分割 $P:\ a = x_0 < x_1 < \cdots < x_n = b$,
  取标志点 $\xi_i \in E_i = [x_{i-1}, x_i]$, 对\alert{积分和}
  $I(f, P) = \sum_{i=1}^{n} f(\xi_i)\, |E_i|$ 取极限 ($\lambda(P) = \max_i |E_i| \to 0$) 来定义;
\item 对某些典型函数证明了其可积性 (极限的存在性), 但仅凭定义\alert{计算积分值}相当困难:
  最终靠 Newton--Leibniz 公式, 在不少情形给出积分的 closed form
  (原函数在两个区间端点之差的差值);
\item \textbf{Lebesgue 积分} (第四章): 对\alert{值域}分割构造积分和, 再取极限
  (\S 4.1 引言中的类比).
\end{itemize}

\pause

\begin{itemize}
\item 本节目标: 联系新旧两种积分理论:
  \begin{enumerate}
  \item L 积分是 R 积分的\alert{推广}: 在正常 R 积分范围内, $R_E \subset L_E$, 且在 $R_E$ 上积分值一致;
  \item 通过比较发现: L 积分理论在大部分时候是更优的一套理论, 能站在更高的角度,
    解释、证明 R 理论\alert{难以处理}的结论, 特别是\alert{极限与积分换序}.
  \end{enumerate}
\end{itemize}

\end{frame}

%------------------------------------------------

\section[两准则]{黎曼可积的两个准则}

%------------------------------------------------

\begin{frame}{回忆: 函数黎曼可积的两个准则}

\begin{attnblock}[黎曼可积准则]
$f$ 在 $E = [a, b]$ 上黎曼可积 (记作 $f \in R_E$)
$\Longleftrightarrow$ $f$ 在 $E$ 上\alert{有界}, 且下面两条之一成立:
\textcircled{1} (达布) 上、下积分相等; \quad
\textcircled{2} (勒贝格) 几乎处处连续.
\end{attnblock}

\pause

\begin{itemize}
\item 记号: $|E_i| = \Delta x_i$ (子区间长度), $\lambda(P) = \max_i |E_i|$ (划分细度),
  $I(f, P) = \sum_{i=1}^{n} f(\xi_i)\, |E_i|$ (带标志点的积分和);
\item 数学分析中把它们作为定理给出; 本节依次重证:
  \alert{1$^\circ$} 可积 $\Rightarrow$ 有界, \quad
  \alert{2$^\circ$} 达布准则, \quad
  \alert{3$^\circ$} 勒贝格准则;
\item 其中 3$^\circ$ 将真正动用勒贝格积分的工具 (有界收敛定理、积分唯一性),
  ``新理论解释旧理论''的第一个实例.
\end{itemize}

\end{frame}

%------------------------------------------------

\section[有界性]{1$^\circ$: 可积蕴含有界}

%------------------------------------------------

\begin{frame}{1$^\circ$: 黎曼可积 $\Longrightarrow$ 有界}

\begin{prop}
$f \in R_E \ \Longrightarrow\ f$ 在 $E$ 上有界.
\end{prop}

\startproof[bg=yellow, fg=red](证明: 反证法)
假设 $f$ 在 $E$ 上无界, 那么对 $E$ 的\alert{任意}划分 $P$, 至少有一个子区间
$E_k = [x_{k-1}, x_k]$, $f$ 在其上无界. 取 $X = 1/|E_k| > 0$,
由无界性 $\exists\, \xi_k \in E_k$, 使得 $|f(\xi_k)| > X$;
再取 $X' = 2|f(\xi_k)|$, 又 $\exists\, \xi_k' \in E_k$, 使得 $|f(\xi_k')| > X'$. 于是
\[
|f(\xi_k') - f(\xi_k)|
\geqslant |f(\xi_k')| - |f(\xi_k)|
> 2|f(\xi_k)| - |f(\xi_k)|
> \frac{1}{|E_k|} .
\]
将 $\xi_k$ 与 $\xi_k'$ 分别作为标志点, \alert{其余标志点一致}, 得同一划分 $P$ 的两个积分和:
\[
\Big| \sum_{i=1}^{n} f(\xi_i')\, |E_i| - \sum_{i=1}^{n} f(\xi_i)\, |E_i| \Big|
= |f(\xi_k') - f(\xi_k)| \cdot |E_k| > 1 .
\]
这表明 $\displaystyle \lim_{\lambda(P) \to 0} \sum_{i=1}^{n} f(\xi_i)\, |E_i|$ 的极限不存在. \qed

\pause

\begin{itemize}
\item 要害: 可积要求积分和关于\alert{一切}标志点的选取一致地收敛.
\end{itemize}

\end{frame}

%------------------------------------------------

\section[达布准则]{2$^\circ$: 达布准则}

%------------------------------------------------

\begin{frame}{2$^\circ$: 达布准则}

\begin{thm}[达布准则]
设 $f$ 有界, 则 $f \in R_E \ \Longleftrightarrow\ $ 达布上、下积分相等.
\end{thm}

\begin{itemize}
\item 达布大 (上) 和: $S(f, P) = \sum_{i=1}^{n(P)} M_i\, |E_i|$,
  $M_i = \sup\limits_{x \in E_i} f(x)$; \quad
  达布小 (下) 和: $s(f, P) = \sum_{i=1}^{n(P)} m_i\, |E_i|$,
  $m_i = \inf\limits_{x \in E_i} f(x)$;
\item 上积分与下积分:
$\displaystyle \overline{\int}_a^b f(x) ~ \mathrm{d} x := \inf_P S(f, P)$,
$\displaystyle \underline{\int}_a^b f(x) ~ \mathrm{d} x := \sup_P s(f, P)$.
\end{itemize}

\pause

对一般有界 $f$, \alert{达布定理}保证两个极限都存在:
\[
\underbrace{\lim_{\lambda(P) \to 0} s(f, P)}_{\text{自动存在}}
= \sup_P s(f, P)
= \underline{\int}_a^b f ~ \mathrm{d} x
\;\leqslant\;
\overline{\int}_a^b f ~ \mathrm{d} x
= \inf_P S(f, P)
= \underbrace{\lim_{\lambda(P) \to 0} S(f, P)}_{\text{自动存在}} ,
\]
而 $f \in R_E \Longleftrightarrow$ 中间的 ``$\leqslant$'' 取 ``$=$''.

\end{frame}

%------------------------------------------------

\begin{frame}{达布准则的证明 (``$\Rightarrow$'')}

设 $f \in R_E$, 记 $I = \displaystyle \int_a^b f(x) ~ \mathrm{d} x
= \lim_{\lambda(P) \to 0} \sum_{i=1}^{n} f(\xi_i)\, |E_i|$.

\pause

则 $\forall\, \varepsilon > 0$, $\exists\, \delta > 0$, 使得对一切带标志点的划分 $P$,
只要 $\lambda(P) < \delta$, 就有
\[
\Big| \sum_{i=1}^{n} f(\xi_i)\, |E_i| - I \Big| = |I(f, P) - I| < \frac{\varepsilon}{2} .
\]

由 $M_i = \sup\limits_{x \in E_i} f(x)$, $\exists\, \xi_i \in E_i$ 使得
$|f(\xi_i) - M_i| < \varepsilon / 2|E|$;
由 $m_i = \inf\limits_{x \in E_i} f(x)$, $\exists\, \xi_i' \in E_i$ 使得
$|f(\xi_i') - m_i| < \varepsilon / 2|E|$
(\alert{注意: 这里使用的是带标志点的划分}). 于是
\[
|I(f, P) - S(f, P)|
= \sum_{i=1}^{n(P)} |f(\xi_i) - M_i| \cdot |E_i|
< \frac{\varepsilon}{2|E|} \sum_{i=1}^{n(P)} |E_i| = \frac{\varepsilon}{2},
\qquad
|I(f, P) - s(f, P)| < \frac{\varepsilon}{2} .
\]

\pause

即 $I(f, P) - \dfrac{\varepsilon}{2} < s(f, P) \leqslant S(f, P) < I(f, P) + \dfrac{\varepsilon}{2}$
(凡 $\lambda(P) < \delta$). 令 $\lambda(P) \to 0$, 用达布定理:
\[
I - \frac{\varepsilon}{2}
\;\leqslant\;
\underline{\int}_a^b f ~ \mathrm{d} x
\;\leqslant\;
\overline{\int}_a^b f ~ \mathrm{d} x
\;\leqslant\;
I + \frac{\varepsilon}{2} .
\]
由 $\varepsilon$ 任意, 上、下积分相等.
\qed

\end{frame}

%------------------------------------------------

\begin{frame}{达布准则的证明 (``$\Leftarrow$'')}

设上、下积分相等, 记公共值为 $I$.

对任意带标志点的划分 $P$, 恒有
\[
s(f, P) \;\leqslant\; I(f, P) \;\leqslant\; S(f, P) .
\]

\pause

由达布定理,
$\displaystyle \lim_{\lambda(P) \to 0} s(f, P) = \lim_{\lambda(P) \to 0} S(f, P) = I$,
由夹逼得 $\displaystyle \lim_{\lambda(P) \to 0} I(f, P)$ 极限存在且为 $I$, 即 $f \in R_E$.
\qed

\pause

\begin{itemize}
\item 达布准则的价值: 把``积分和的极限存在'' (难以直接验证)
  转化为``上、下积分相等'' (上、下确界总是存在).
\end{itemize}

\end{frame}

%------------------------------------------------

\section[勒贝格准则]{3$^\circ$: 勒贝格准则}

%------------------------------------------------

\begin{frame}{3$^\circ$: 勒贝格准则与振幅函数}

\begin{thm}[勒贝格准则]
设 $f$ 有界, 则 $f \in R_E \ \Longleftrightarrow\ f$ 在 $E$ 上\alert{几乎处处连续}.
\end{thm}

引入 (回忆) $f$ 在一点 $x \in E$ 处的\alert{振幅}:
\[
\omega_f(x) := \lim_{\delta \to 0^+} \omega\big( f,\, U_E(x, \delta) \big),
\qquad
\omega\big( f,\, U_E(x, \delta) \big)
= \sup_{x_1, x_2 \in U_E(x, \delta)} |f(x_1) - f(x_2)| ,
\]
其中 $U_E(x, \delta) = E \cap U(x, \delta)$ 为相对邻域.

\pause

\begin{itemize}
\item $\omega(f, U_E(x, \delta))$ 随 $\delta \downarrow 0$ 单调不增、有下界, 故极限存在;
  $\omega_f$ 是 $E$ 上的\alert{有界可测}函数: $0 \leqslant \omega_f \leqslant \omega(f, E) \leqslant 2M$
  (可测性见下页);
\item $f$ 在 $x$ 处连续 $\Longleftrightarrow \omega_f(x) = 0$;
  \alert{不连续点集} $= E(\omega_f > 0)$.
\end{itemize}

\end{frame}

%------------------------------------------------

\begin{frame}{Claim: $E(\omega_f < \alpha)$ 是开集}

\begin{lem}[Claim]
$\forall\, \alpha > 0$, $E(\omega_f < \alpha)$ 是 (相对 $E$) 开集.
\end{lem}

\startproof[bg=yellow, fg=red](证明)
$\forall\, x \in E(\omega_f < \alpha)$: 由
$\alpha > \omega_f(x) = \lim\limits_{\delta \to 0^+} \omega(f, U_E(x, \delta))$,
$\exists\, \delta > 0$ (这里的 $\delta$ 与 $x$ 有关), 使得
$\omega(f, U_E(x, \delta)) < \alpha$.

任取 $x' \in U_E(x, \delta)$, 令 $\delta' = (\delta - |x - x'|)/3 > 0$,
则 $U_E(x', \delta') \subset U_E(x, \delta)$, 从而
\[
\omega_f(x') \;\leqslant\; \omega\big( f,\, U_E(x', \delta') \big)
\;\leqslant\; \omega\big( f,\, U_E(x, \delta) \big) \;<\; \alpha ,
\]
即 $x' \in E(\omega_f < \alpha)$. 由 $x'$ 任取自 $U_E(x, \delta)$,
故 $U_E(x, \delta) \subset E(\omega_f < \alpha)$, $x$ 是\alert{内点}.

以上即证明了 $E(\omega_f < \alpha)$ 的所有点都是内点, 从而 $E(\omega_f < \alpha)$ 为开集.
\qed

\pause

\begin{itemize}
\item 由此, $E(\omega_f < \alpha)$ 可测 ($\forall \alpha$), 即 $\omega_f$ 是可测函数.
\end{itemize}

\end{frame}

%------------------------------------------------

\begin{frame}{命题: 振幅的积分 $=$ 上积分 $-$ 下积分}

\begin{prop}
$\displaystyle \int_E \omega_f ~ \mathrm{d} m
= \overline{\int}_a^b f(x) ~ \mathrm{d} x - \underline{\int}_a^b f(x) ~ \mathrm{d} x$.
\end{prop}

\startproof[bg=yellow, fg=red](证明)
取一列划分 $\{P_k\}_{k \in \mathbb{N}}$, $\lambda(P_k) \to 0$, 令 (非负简单函数)
\[
\omega_{P_k}(x) =
\begin{cases}
M_i^{(k)} - m_i^{(k)} , & x \in \big( x_{i-1}^{(k)},\, x_i^{(k)} \big), \\
0 , & x \in \big\{ x_0^{(k)},\, x_1^{(k)},\, \dots,\, x_{n_k}^{(k)} \big\} ;
\end{cases}
\]
记 $A = \bigcup_{k \geqslant 1} A(P_k)$ (全部分点, 可数), $mA = 0$;
且 $\forall\, x \in E \setminus A$, $\lim\limits_{k} \omega_{P_k}(x) = \omega_f(x)$
(小区间长度 $\leqslant \lambda(P_k) \to 0$).
\qed

\end{frame}

%------------------------------------------------

\begin{frame}{命题证明 (续): 有界收敛定理}

\startproof[bg=yellow, fg=red](证明 (续))
设 $|f| \leqslant M$, 则 $0 \leqslant \omega_{P_k}(x) \leqslant \omega(f, E) \leqslant 2M$,
$mE = b - a < \infty$; 由\alert{有界收敛定理} (末步极限由达布定理保证存在):
\[
\begin{aligned}
\int_E \omega_f ~ \mathrm{d} m
&= \int_E \lim_{k \to \infty} \omega_{P_k} ~ \mathrm{d} m
= \lim_{k \to \infty} \int_E \omega_{P_k} ~ \mathrm{d} m \\
&= \lim_{k \to \infty} \Big( \sum_{i=1}^{n_k} M_i^{(k)} |E_i^{(k)}| - \sum_{i=1}^{n_k} m_i^{(k)} |E_i^{(k)}| \Big)
= \lim_{k \to \infty} S(f, P_k) - \lim_{k \to \infty} s(f, P_k) \\
&= \overline{\int}_a^b f(x) ~ \mathrm{d} x - \underline{\int}_a^b f(x) ~ \mathrm{d} x .
\end{aligned}
\]
\qed

\pause

\begin{itemize}
\item 这正是勒贝格积分的用武之地: 达布定理 $+$ 有界收敛定理
  把``上、下积分之差''化成一个\alert{可测函数的积分}.
\end{itemize}

\end{frame}

%------------------------------------------------

\begin{frame}{勒贝格准则的证明}

由达布准则, $f \in R_E \Longleftrightarrow \overline{\int}_a^b f ~ \mathrm{d} x = \underline{\int}_a^b f ~ \mathrm{d} x$,
由上面的命题, 这等价于
\[
\int_E \omega_f ~ \mathrm{d} m = 0 .
\]

\pause

由勒贝格积分\alert{唯一性定理} (\S 4.2: $\int_E |h| ~ \mathrm{d} m = 0 \Longleftrightarrow h \sim 0$,
这里 $\omega_f \geqslant 0$):
\[
\int_E \omega_f ~ \mathrm{d} m = 0
\;\Longleftrightarrow\;
\omega_f \sim 0
\;\Longleftrightarrow\;
mE(\omega_f > 0) = 0 .
\]

\pause

而 $f$ 的不连续点集 $= E(\omega_f > 0)$, 故以上这一串等价关系即
\[
f \in R_E
\;\Longleftrightarrow\;
f \text{ 在 } E \text{ 上几乎处处连续} .
\]
\qed

\pause

\begin{itemize}
\item 证明真正用到了勒贝格积分 (有界收敛定理 $+$ 唯一性定理),
  站在更高的角度看旧理论.
\end{itemize}

\end{frame}

%------------------------------------------------

\section[R$\subset$L]{从 R 积分到 L 积分}

%------------------------------------------------

\begin{frame}{定理: 黎曼可积 $\Longrightarrow$ 勒贝格可积, 且积分相等}

\begin{thm}
设 $E = [a, b]$, 若 $f \in R_E$, 则 $f \in L_E$, 且积分值相等.
\end{thm}

\startproof[bg=yellow, fg=red](证明)
由勒贝格准则, $f \in R_E \Rightarrow f$ 有界且 a.e.\ 连续
$\Rightarrow f$ 可测 (零测集外连续) $\Rightarrow f \in L_E$ (\S 4.1: $mE < \infty$ 且有界).

对任意划分 $P$ (子区间 $E_i = [x_{i-1}, x_i]$):
$m_i |E_i| = \int_{E_i} m_i ~ \mathrm{d} m \leqslant \int_{E_i} f ~ \mathrm{d} m \leqslant \int_{E_i} M_i ~ \mathrm{d} m = M_i |E_i|$,
求和得 $s(f, P) \leqslant \int_E f ~ \mathrm{d} m \leqslant S(f, P)$; 取 $\sup_P$, $\inf_P$:
\[
\underline{\int}_a^b f(x) ~ \mathrm{d} x = \sup_P s(f, P)
\;\leqslant\; \int_E f ~ \mathrm{d} m
\;\leqslant\; \inf_P S(f, P) = \overline{\int}_a^b f(x) ~ \mathrm{d} x .
\]
由达布准则 ($f \in R_E$: 上 $=$ 下积分), 上式 ``$\leqslant$'' 全取 ``$=$'', 即
\[
\int_E f ~ \mathrm{d} m
= \underline{\int}_a^b f(x) ~ \mathrm{d} x
= \overline{\int}_a^b f(x) ~ \mathrm{d} x
= (R) \int_a^b f(x) ~ \mathrm{d} x .
\]
\qed

\end{frame}

%------------------------------------------------

\begin{frame}{注: $L_{[a,b]}$ 是 $R_{[a,b]}$ 的``完备化''}

\begin{itemize}
\item 结合计较完备的事实: $E$ 上可测函数空间关于\alert{依测度收敛}是完备的 (\S 3.2), 我们可以说
  \[
  L_{[a,b]} \big/ \sim
  \;\supset\;
  R_{[a,b]} \big/ \sim :
  \quad
  L_{[a,b]} \text{ 是 } R_{[a,b]} \text{ 的 ``完备化''} ;
  \]
\graymode
\item 具体地: $f_n \in R_{[a,b]}$ a.e.\ (或依测度) 收敛于 $f$ 且有可积控制, 则 $f$ 未必黎曼可积,
  但必有 $f \in L_{[a,b]}$. 例如 $[0,1]$ 上的 Dirichlet 函数
  $\chi_{\mathbb{Q} \cap [0,1]}$: 取 $\{q_k\} = \mathbb{Q} \cap [0, 1]$ 的枚举,
  阶梯函数 $f_n = \chi_{\{q_1, \dots, q_n\}} \in R_{[0,1]}$ 渐升地
  $f_n \uparrow \chi_{\mathbb{Q}} \sim 0 \in L_{[0,1]}$,
  而 $\chi_{\mathbb{Q} \cap [0,1]} \notin R_{[0,1]}$ (处处不连续);
\normalmode
\item 下一页的黎曼函数是更精致的例子: $f \in R_{[0,1]}$, 但其``性质''只有用测度语言才看得清.
\end{itemize}

\end{frame}

%------------------------------------------------

\section[黎曼函数]{例: 黎曼函数}

%------------------------------------------------

\begin{frame}{例: 黎曼函数 (定义与不连续点)}

\begin{eg}
$E = [0, 1]$ 上的黎曼函数
\[
f(x) =
\begin{cases}
1/q , & x = p/q \text{ 为既约分数 } (q \geqslant 1), \\
0 , & x \in [0,1] \setminus \mathbb{Q} \ \cup\ \{0\} .
\end{cases}
\]
\end{eg}

\startproof[bg=yellow, fg=red]($1^\circ$ $f$ 在每个既约分数点不连续)
设 $x = p/q \in (0,1) \cap \mathbb{Q}$ 为既约分数. $\forall\, \delta > 0$,
$U(x, \delta) \cap \big( [0,1] \setminus \mathbb{Q} \big) \neq \varnothing$
(无理点稠密), 邻域内振幅至少跨过 $1/q$ 与 $0$:
\[
\omega_f(x) \;\geqslant\; |1/q - 0| = 1/q \;>\; 0 .
\]
\qed

\pause

\begin{itemize}
\item 不连续点集 $(0,1) \cap \mathbb{Q}$ \alert{稠密}于 $[0,1]$,
  经典意义下``性质很差''的函数.
\end{itemize}

\end{frame}

%------------------------------------------------

\begin{frame}{例: 黎曼函数 (连续点与积分值)}

\startproof[bg=yellow, fg=red]($2^\circ$ $f$ 在每个无理点连续)
事实上, $\forall\, \varepsilon > 0$, 取 $q_0 \in \mathbb{N}$ 使得 $1/q_0 < \varepsilon$;
而
\[
A = \big\{ x = p/q \in (0,1) \cap \mathbb{Q}:\ \text{既约},\ q < q_0 \big\}
\]
为\alert{有限集}. 于是 $\forall\, x \in (0,1) \setminus \mathbb{Q}$, $x \notin A$,
取 $\delta = \rho(x, A)/2 > 0$, 即有 $U(x, \delta) \cap A = \varnothing$.
这意味着 $\forall\, x' \in U(x, \delta)$: 或者 $x'$ 为无理数, $f(x') = 0$;
或者 $x' = p/q$ 既约且 $q \geqslant q_0$, $f(x') = 1/q \leqslant 1/q_0 < \varepsilon$. 所以
\[
\omega_f(x) \leqslant \omega\big( f,\, U(x, \delta) \big)
:= \sup_{x_1, x_2 \in U(x, \delta)} |f(x_1) - f(x_2)| < \varepsilon .
\]
由 $\varepsilon > 0$ 任意, $\omega_f(x) = 0$, 即 $f$ 在 $x$ 处连续. \qed

\pause

由 $1^\circ, 2^\circ$, $f$ 恰在 $(0,1) \cap \mathbb{Q}$ 上不连续, $mE(\text{不连续点}) = 0$,
故 $f$ 在 $E$ 上几乎处处连续, $f \in R_E$, 从而 $f \in L_E$, 且
\[
(R) \int_0^1 f(x) ~ \mathrm{d} x
= \int_E f ~ \mathrm{d} m
= \Big( \int_{E \setminus \mathbb{Q}} + \int_{E \cap \mathbb{Q}} \Big) f ~ \mathrm{d} m
= \int_{E \setminus \mathbb{Q}} 0 ~ \mathrm{d} m + \int_{E \cap \mathbb{Q}} f ~ \mathrm{d} m
= 0 .
\]

\end{frame}

%------------------------------------------------

\section[渐张集列]{渐张集列上的积分}

%------------------------------------------------

\begin{frame}{定理: 渐张集列上的积分 (可积性)}

\begin{thm}
设 $\{E_n\}$ 是渐张可测集列: $E_1 \subset E_2 \subset \cdots \subset E_n \subset \cdots$,
$E = \bigcup_{n=1}^{\infty} E_n$. 设 $f \in L_{E_n}$, $n \in \mathbb{N}$.
那么, 若极限 $\displaystyle \lim_{n \to \infty} \int_{E_n} |f| ~ \mathrm{d} m$ 存在且有限,
则 $f \in L_E$, 且 $\displaystyle \int_E f ~ \mathrm{d} m = \lim_{n \to \infty} \int_{E_n} f ~ \mathrm{d} m$.
\end{thm}

\startproof[bg=yellow, fg=red](证明: 第一步, $f \in L_E$)
考虑 $\big\{ |f| \cdot \chi_{E_n} \big\}_{n \in \mathbb{N}}$, 它是一个\alert{非负渐升}可测函数列,
并且 $\forall\, x \in E$,
$\displaystyle \lim_{n \to \infty} |f(x)| \cdot \chi_{E_n}(x) = |f(x)|$.
由 \alert{Levi 定理}:
\[
\int_E |f| ~ \mathrm{d} m
= \int_E \Big( \lim_{n \to \infty} |f| \cdot \chi_{E_n} \Big) \mathrm{d} m
= \lim_{n \to \infty} \int_E |f| \cdot \chi_{E_n} ~ \mathrm{d} m
= \lim_{n \to \infty} \int_{E_n} |f| ~ \mathrm{d} m
< \infty
\ \Longrightarrow\ f \in L_E .
\]
\qed

\pause

\begin{itemize}
\item Levi 管 $|f|$ (定\alert{可积性}), 控制收敛管 $f$ (求\alert{积分值}),
  见下页.
\end{itemize}

\end{frame}

%------------------------------------------------

\begin{frame}{定理证明 (续): 极限与积分换序}

\startproof[bg=yellow, fg=red](第二步: 求积分值)
又由于 $\forall\, x \in E$,
$\displaystyle \lim_{n \to \infty} f(x)\, \chi_{E_n}(x) = f(x)$,
且 $|f \cdot \chi_{E_n}| \leqslant |f| \in L_E$ ($\forall\, n \in \mathbb{N}$),
由勒贝格\alert{控制收敛定理}:
\[
\int_E f ~ \mathrm{d} m
= \int_E \Big( \lim_{n \to \infty} f \cdot \chi_{E_n} \Big) \mathrm{d} m
= \lim_{n \to \infty} \int_E f \cdot \chi_{E_n} ~ \mathrm{d} m
= \lim_{n \to \infty} \int_{E_n} f ~ \mathrm{d} m .
\]
\qed

\pause

\begin{itemize}
\item 这是 \S 4.3 换序引擎的直接应用, 也是计算\alert{无界区域}上勒贝格积分的标准通道
  (见接下来注 (2));
\item 注意条件是 $\lim_n \int_{E_n} |f| ~ \mathrm{d} m < \infty$: 勒贝格积分天生只认
  \alert{绝对收敛}.
\end{itemize}

\end{frame}

%------------------------------------------------

\section[广义积分]{注: 与广义积分的比较}

%------------------------------------------------

\begin{frame}{注 (1): L 积分是绝对收敛的广义积分}

以上定理表明: 在广义积分对应的情形下, \alert{勒贝格积分指的是绝对收敛的积分},
而广义积分\alert{不要求}绝对收敛:
\[
\int_a^{+\infty} f(x) ~ \mathrm{d} x := \lim_{b \to +\infty} \int_a^b f(x) ~ \mathrm{d} x,
\qquad
\int_a^b f(x) ~ \mathrm{d} x := \lim_{\varepsilon \to 0^+} \int_a^{b - \varepsilon} f(x) ~ \mathrm{d} x .
\]

\pause

\begin{eg}
$\displaystyle \int_0^{+\infty} \frac{\sin x}{x} ~ \mathrm{d} x = \frac{\pi}{2}$ (条件收敛), 但
\[
\int_0^{+\infty} \Big| \frac{\sin x}{x} \Big| ~ \mathrm{d} x
\;\geqslant\; \sum_{k=0}^{\infty} \int_{k\pi + \pi/4}^{k\pi + 3\pi/4} \frac{|\sin x|}{x}\, \mathrm{d} x
\;\geqslant\; \sum_{k=0}^{\infty} \frac{\sqrt{2}/2 \cdot \pi/2}{(k+1)\pi}
= +\infty
\]
(该区间上 $|\sin x| \geqslant \sqrt{2}/2$, $x \leqslant (k+1)\pi$; 调和级数发散).
\end{eg}

故 $f(x) = \dfrac{\sin x}{x} \notin L_{[0, +\infty)}$
(由于 $|f| \notin L_{[0, +\infty)}$).

\end{frame}

%------------------------------------------------

\begin{frame}{注 (2): 无界区域上 L 积分的计算}

由渐张集列定理: $\displaystyle \int_E f ~ \mathrm{d} m = \lim_{n \to \infty} \int_{E_n} f ~ \mathrm{d} m$
可计算某些 $E$ \alert{无界}时的勒贝格积分值.

\begin{itemize}
\item 例如 $E = [0, +\infty)$, 取 $E_n = [0, n]$, 只要 $f \in R_{E_n}$, 则
  $\displaystyle \int_E f ~ \mathrm{d} m = \lim_{n \to \infty} (R) \int_0^n f(x) ~ \mathrm{d} x$;
\item \alert{需注意的是}, 该极限与 $\{E_n\}$ 的选取无关, 只要 $E_n$ 渐张 $\to E$;
\graymode
\item 例: $\displaystyle \int_0^{+\infty} e^{-x^2} ~ \mathrm{d} x
= \lim_{n \to \infty} (R) \int_0^n e^{-x^2} ~ \mathrm{d} x = \frac{\sqrt{\pi}}{2}$
  (概率积分, 每个 $\int_0^n$ 都是正常 Riemann 积分).
\normalmode
\end{itemize}

\end{frame}

%------------------------------------------------

\begin{frame}{注 (3): 反常积分不具备可数可加性}

勒贝格积分对积分区域有\alert{可数可加性}: 若 $E = \bigsqcup_{n=1}^{\infty} E_n$
为互不相交的可测集之并, 则 $\int_E f ~ \mathrm{d} m = \sum_{n=1}^{\infty} \int_{E_n} f ~ \mathrm{d} m$.
但\alert{反常积分可能不具备}这种性质:

\pause

\begin{eg}
$E_n = [n-1, n)$, $E = [0, +\infty)$,
$f = \sum_{n=1}^{\infty} \dfrac{(-1)^{n-1}}{n}\, \chi_{E_n}$. 那么
\[
(R) \int_0^{+\infty} f ~ \mathrm{d} x = \sum_{n=1}^{\infty} \frac{(-1)^{n-1}}{n} = \ln 2
\]
($\ln(1+x)$ 在 $x = 1$ 处的 Taylor 级数). 但该级数仅\alert{条件收敛}:
由 Riemann 重排定理, $\forall\, -\infty \leqslant \alpha \leqslant \beta \leqslant +\infty$,
存在重排 $\tau_{\alpha, \beta}$, 令 $S_k^{\alpha, \beta} = \sum_{n=1}^{k} a_{\tau_{\alpha,\beta}(n)}$, 有
$\varliminf\limits_{k \to \infty} S_k^{\alpha,\beta} = \alpha$,
$\varlimsup\limits_{k \to \infty} S_k^{\alpha,\beta} = \beta$.

特别地取 $\alpha = \beta \neq \ln 2$, 以及 $E_n' = E_{\tau_{\alpha,\beta}(n)}$, 那么
$E = \bigsqcup_{n=1}^{\infty} E_n'$ 也是互不相交的可测集之并, 且
\[
\sum_{n=1}^{\infty} \int_{E_n'} f ~ \mathrm{d} x
= \lim_{k \to \infty} S_k^{\alpha, \beta} = \alpha \;\neq\; \ln 2 = \int_E f ~ \mathrm{d} x .
\]
\end{eg}

黎曼积分的两个推广: $\longrightarrow$ 广义积分 (对积分区域无可数可加性);
$\longrightarrow$ 勒贝格积分 (绝对可积).

\end{frame}

%------------------------------------------------

\section[Fourier]{Fourier 级数的逐项积分}

%------------------------------------------------

\begin{frame}{定理: Fourier 级数的逐项积分}

\begin{thm}[逐项积分]
设 $f \in L_{[-\pi, \pi]}$,
$f(x) \sim \dfrac{a_0}{2} + \sum_{n=1}^{\infty} \big( a_n \cos nx + b_n \sin nx \big)$,
则 $\forall\, \alpha, \beta \in [-\pi, \pi]$, 有
\[
\int_{\alpha}^{\beta} f(x) ~ \mathrm{d} x
= \int_{\alpha}^{\beta} \frac{a_0}{2}\, \mathrm{d} x
+ \sum_{n=1}^{\infty} \int_{\alpha}^{\beta}
\big( a_n \cos nx + b_n \sin nx \big) \mathrm{d} x .
\]
\end{thm}

\pause

\begin{itemize}
\item 无需 Fourier 级数\alert{收敛}于 $f$, 甚至无需级数点点收敛: 逐项积分\alert{恒}成立;
\item 这是勒贝格积分优于 Riemann 积分的标志性结果: R 积分框架下需要对级数的一致收敛性
  作额外假定.
\end{itemize}

\end{frame}

%------------------------------------------------

\begin{frame}{证明 (一): Fourier 部分和逼近}

\startproof[bg=yellow, fg=red](证明)
设 $\varphi = \chi_{(\alpha, \beta)}$ 为定义在 $[-\pi, \pi]$ 上的函数,并将其延拓至
$(-\infty, +\infty)$ 上, 使其有周期 $2\pi$. 那么 $\varphi$ 有 Fourier 展式
\[
\varphi(x) \sim \frac{A_0}{2} + \sum_{n=1}^{\infty}
\big( A_n \cos nx + B_n \sin nx \big),
\]
其中
\[
A_n = \frac{1}{\pi} \int_{\alpha}^{\beta} \cos nx ~ \mathrm{d} x
= \frac{\sin n\beta - \sin n\alpha}{n\pi},
\qquad
B_n = \frac{1}{\pi} \int_{\alpha}^{\beta} \sin nx ~ \mathrm{d} x
= \frac{\cos n\alpha - \cos n\beta}{n\pi} .
\]
$\varphi$ 分段光滑, 除 $x = \alpha, \beta$ 外处处等于其 Fourier 级数的和
(在 $\alpha, \beta$ 处级数收敛于左右平均值), 故令
$\varphi_n = \dfrac{A_0}{2} + \sum_{k=1}^{n} (A_k \cos kx + B_k \sin kx)$, 有
$\lim\limits_{n \to \infty} \varphi_n(x) = \varphi(x)$ 对 a.e.\ $x \in \mathbb{R}$.
\qed

\end{frame}

%------------------------------------------------

\begin{frame}{证明 (二): 逐项比较}

\startproof[bg=yellow, fg=red](证明 (续))
现在, $\int_{\alpha}^{\beta} f(x) ~ \mathrm{d} x = \int_{-\pi}^{\pi} f(x)\, \varphi(x) ~ \mathrm{d} x$;
另一方面, 由 $a_n = \frac{1}{\pi} \int_{-\pi}^{\pi} f \cos nx ~ \mathrm{d} x$ 等,
\[
\int_{\alpha}^{\beta} a_n \cos nx ~ \mathrm{d} x
= a_n \cdot \frac{\sin n\beta - \sin n\alpha}{n}
= A_n \int_{-\pi}^{\pi} f(x) \cos nx ~ \mathrm{d} x ,
\]
同理 $\int_{\alpha}^{\beta} b_n \sin nx ~ \mathrm{d} x
= B_n \int_{-\pi}^{\pi} f(x) \sin nx ~ \mathrm{d} x$. 从而对每个 $n$,
\[
\int_{\alpha}^{\beta} \frac{a_0}{2}\, \mathrm{d} x
+ \sum_{k=1}^{n} \int_{\alpha}^{\beta}
\big( a_k \cos kx + b_k \sin kx \big) \mathrm{d} x
= \int_{-\pi}^{\pi} \varphi_n(x) \cdot f(x) ~ \mathrm{d} x .
\]
\qed

\end{frame}

%------------------------------------------------

\begin{frame}{证明 (三): 问题与策略}

于是, 定理转化为了对
\[
\int_{\alpha}^{\beta} f(x) ~ \mathrm{d} x
= \lim_{n \to \infty} \int_{-\pi}^{\pi} f(x) \cdot \varphi_n(x) ~ \mathrm{d} x
\]
的证明.

\pause

\begin{itemize}
\item 已有 $\varphi_n(x) \to \varphi(x)$ a.e., 自然想用\alert{控制收敛定理};
\item \textbf{※} 若有 $g \in L_{[-\pi, \pi]}$, 使得
  $|f(x)\, \varphi_n(x)| \leqslant g(x)$ 对 a.e.\ $x \in [-\pi, \pi]$ 成立,
  那么上式即可用控制收敛定理得到;
\item 因此只需: $\varphi_n$ \alert{一致有界}, 即 $\exists\, C \in \mathbb{R}_+$, 使得
  $|\varphi_n(x)| \leqslant C$ ($\forall\, n \in \mathbb{N}$, $x \in [-\pi, \pi]$);
  那么 $g = C|f|$ 就是想要的控制函数.
\end{itemize}

\end{frame}

%------------------------------------------------

\begin{frame}{证明 (四): Dirichlet 核表示}

\startproof[bg=yellow, fg=red](Claim $\varphi_n$ 一致有界: 核表示)
实际上, 由 Dirichlet 核 ($\varphi = \chi_{(\alpha,\beta)}$)
\[
\varphi_n(x) = \frac{1}{\pi} \int_{-\pi}^{\pi} \varphi(x + u)\,
\frac{\sin \dfrac{2n+1}{2}\, u}{2 \sin \dfrac{u}{2}} ~ \mathrm{d} u
= \frac{1}{\pi} \int_{\alpha - x}^{\beta - x} D_n(u) ~ \mathrm{d} u
= \frac{2}{\pi} \int_{\frac{\alpha - x}{2}}^{\frac{\beta - x}{2}} D_n(2v) ~ \mathrm{d} v ,
\]
其中 $D_n(u) = \dfrac{\sin \tfrac{2n+1}{2} u}{2 \sin \tfrac{u}{2}}$ (Dirichlet 核).

由于 $D_n(-u) = D_n(u)$ (偶), 且 $D_n(2v)$ 关于 $v$ 以 $\pi$ 为周期,
只需考虑 $0 \leqslant \alpha - x < \beta - x \leqslant \pi$ 的情形.
\qed

\pause

\begin{center}
\begin{tikzpicture}[scale = 0.94]
\draw[->] (-0.3, 0) -- (8.6, 0) node[right] {$v$};
\draw (0, 0.1) -- (0, -0.1) node[below] {$0$};
\draw (2.4, 0.08) -- (2.4, -0.08) node[below] {$\frac{\pi}{2n+1}$};
\draw (4.4, 0.08) -- (4.4, -0.08) node[below] {$\frac{2\pi}{2n+1}$};
\draw (6.4, 0.08) -- (6.4, -0.08) node[below] {$\frac{3\pi}{2n+1}$};
\draw (8.0, 0.08) -- (8.0, -0.08) node[below] {$\frac{\pi}{2}$};
\node at (1.2, 0.4) {$I_0$};
\node at (3.4, 0.4) {$I_1$};
\node at (5.4, 0.4) {$I_2$};
\node at (7.3, 0.4) {$\cdots$};
\node[themecolor] at (1.2, -0.65) {$+$};
\node[themecolor] at (3.4, -0.65) {$-$};
\node[themecolor] at (5.4, -0.65) {$+$};
\node[themecolor] at (7.3, -0.65) {$-$};
\end{tikzpicture}
\end{center}

$I_k = \Big[ \dfrac{k\pi}{2n+1},\ \dfrac{(k+1)\pi}{2n+1} \Big]$:
其上 $\sin (2n+1)v$ 完成一个拱形, 符号 $(-1)^k$ 交替, 而 $|\sin v|$ 逐段抬高.

\end{frame}

%------------------------------------------------

\begin{frame}{证明 (五): 交错区间估计}

\startproof[bg=yellow, fg=red](Claim $\varphi_n$ 一致有界 (续))
在相邻区间 $I_k$, $I_{k+1}$ 上有
$T_k := \int_{I_k} D_n(2v) ~ \mathrm{d} v$ 满足
\[
|T_{k+1}| < |T_k|,
\qquad
T_k \text{ 与 } T_{k+1} \text{ 符号相反} .
\]
把 $\big[ \tfrac{\alpha - x}{2}, \tfrac{\beta - x}{2} \big]$ 上的积分写成
``左端块 $+$ 若干完整 $I_k$ $+$ 右端块'': 端块不超过相邻完整段, 而完整段两两相消后
绝对值不超过首个完整段. 因此, 无论首个完整段符号如何,
\[
|\varphi_n(x)|
\;\leqslant\; 2 \cdot \frac{2}{\pi} \cdot \max_k \int_{I_k} |D_n(2v)| ~ \mathrm{d} v .
\]
\qed

\pause

由 $|\sin v|$ 递增, $\max_k$ 在 $I_0$ 取到. 记 $L = \dfrac{\pi}{2n+1}$:
\[
\int_{I_0} |D_n(2v)| ~ \mathrm{d} v
= \int_0^{L} \frac{\sin (2n+1)v}{\sin v} ~ \mathrm{d} v
\;\leqslant\; \int_0^{L} \frac{(2n+1)v}{\sin v} ~ \mathrm{d} v
\;\leqslant\; (2n+1) \cdot L \cdot \frac{L}{\sin L}
\]
(用 $\sin t \leqslant t$ 及 $\dfrac{v}{\sin v}$ 在 $(0, \pi)$ 单调增). 于是
\[
|\varphi_n(x)| \;\leqslant\; \frac{4}{\pi} \cdot (2n+1)\, L \cdot \frac{L}{\sin L}
= \frac{4}{\pi} \cdot \pi \cdot \frac{L}{\sin L}
\;\leqslant\; \frac{4}{\pi} \cdot \pi \cdot \frac{\pi/3}{\sin(\pi/3)}
= \frac{8\pi}{3\sqrt{3}}
\;<\; 2\pi ,
\]
与 $n$, $x$ 都无关: $|\varphi_n(x)| \leqslant C := 2\pi$.

\end{frame}

%------------------------------------------------

\begin{frame}{证明 (六): 控制收敛收尾}

现在 $f \cdot \varphi_n \to f \cdot \varphi$ a.e.\ 于 $[-\pi, \pi]$, 且
\[
|f(x)\, \varphi_n(x)| \;\leqslant\; C\, |f(x)| \;\in\; L_{[-\pi, \pi]} ,
\]
由\alert{控制收敛定理}:
\[
\lim_{n \to \infty} \int_{-\pi}^{\pi} f(x)\, \varphi_n(x) ~ \mathrm{d} x
= \int_{-\pi}^{\pi} f(x)\, \varphi(x) ~ \mathrm{d} x
= \int_{\alpha}^{\beta} f(x) ~ \mathrm{d} x .
\]
代回证明 (二) 末的恒等式:
\[
\int_{\alpha}^{\beta} f(x) ~ \mathrm{d} x
= \lim_{n \to \infty} \Big[ \int_{\alpha}^{\beta} \frac{a_0}{2}\, \mathrm{d} x
+ \sum_{k=1}^{n} \int_{\alpha}^{\beta}
\big( a_k \cos kx + b_k \sin kx \big) \mathrm{d} x \Big]
\]
\[
= \int_{\alpha}^{\beta} \frac{a_0}{2}\, \mathrm{d} x
+ \sum_{k=1}^{\infty} \int_{\alpha}^{\beta}
\big( a_k \cos kx + b_k \sin kx \big) \mathrm{d} x .
\]
\qed

\pause

\begin{itemize}
\item 回顾: 难点全部集中在 $|\varphi_n|$ 的\alert{一致有界性}:
  用 Dirichlet 核的交错结构硬估计出来; 换序本身则由控制收敛定理一步完成.
\end{itemize}

\end{frame}

%------------------------------------------------

\section[小结]{小结}

%------------------------------------------------

\begin{frame}{小结: 本节结论链}

\begin{itemize}
\item \textbf{三个准则的重证}: 1$^\circ$ 可积 $\Rightarrow$ 有界 (反证);
  2$^\circ$ 达布准则 (上、下积分);
  3$^\circ$ 勒贝格准则 (振幅 $\omega_f$ $+$ 有界收敛 $+$ 唯一性):
  用 L 工具解释 R 理论;
\item \textbf{推广定理}: $R_E \subset L_E$ 且积分值相等;
  $L_{[a,b]}$ 是 $R_{[a,b]}$ 的``完备化'';
\item \textbf{例}: 黎曼函数 (恰在有理点不连续, a.e.\ 连续, R、L 积分均为 $0$);
  渐张集列 ($E_n \uparrow E$: Levi 定可积 $+$ 控制收敛求值);
\item \textbf{广义积分三注}: L 积分 $=$ 绝对收敛型 ($\sin x / x$ 反例);
  无界区域的计算通道 (与 $\{E_n\}$ 选取无关);
  反常积分无可数可加性 (条件收敛级数的重排, $\ln 2$ 可变任值);
\item \textbf{Fourier 逐项积分}: $\varphi = \chi_{(\alpha,\beta)}$ 的 Fourier 部分和
  $\varphi_n \to \varphi$ a.e.\ 且 $|\varphi_n| \leqslant 2\pi$ 一致有界 (Dirichlet 核交错估计),
  控制收玫定理完成换序.
\end{itemize}

\pause

\begin{itemize}
\item 展望: 下一节 \S 4.5 乘积测度与 Fubini 定理: 高维情形的积分换序;
\item 作业: 教材第四章习题 \textbf{21, 23, 25, 26}.
\end{itemize}

\end{frame}

%------------------------------------------------

\ifIncludeExercise
\begin{frame}{作业}

\begin{block}{教材第四章习题}
\begin{itemize}
\item \textbf{21.} 设 $f$ 在 $(-\infty, +\infty)$ 上可积, 证明
$\displaystyle \lim_{h \to 0} \int_{-\infty}^{+\infty} |f(x+h) - f(x)| ~ \mathrm{d} m = 0$.
\item \textbf{23.} 设 $f$ 是 $\mathbb{R}$ 上的可积函数, 试证
$\hat{f}(t) = \int_{\mathbb{R}} e^{-itx} f(x) ~ \mathrm{d} x$
是 $\mathbb{R}$ 上的连续函数, 且
$\hat{f}(t) = \dfrac{\mathrm{d}}{\mathrm{d} t} \int_{\mathbb{R}} \dfrac{e^{-itx} - 1}{-ix}\, f(x) ~ \mathrm{d} x$.
\end{itemize}
\end{block}

\begin{block}{续}
\begin{itemize}
\item \textbf{25.} 设 $f$ 是 $\mathbb{R}$ 上的可测函数, 令
$\mu(\alpha) = m\mathbb{R}(|f| > \alpha)$, 试证
$\displaystyle \int_{\mathbb{R}} |f|^p ~ \mathrm{d} m
= p \int_0^{+\infty} \alpha^{p-1} \mu(\alpha) ~ \mathrm{d} \alpha$
($1 \leqslant p < \infty$).
\item \textbf{26.} 设 $mE < \infty$, 证明函数 $f$ 在 $E$ 上可积的充分必要条件是级数
$\displaystyle \sum_{n=1}^{\infty} mE(|f| \geqslant n)$ 收敛.
当 $mE = \infty$ 时, 结论是否成立?
\end{itemize}
\end{block}

\end{frame}
\fi

%------------------------------------------------

\end{document}
